\cotefile{162-5}
\scanpage{29}
On peut alors de façon unique prolonger en une \deleted{dérivation} opération $\delta_M$ sur $\bigwedge^*\Omega\otimes M$, en posant
\[
\delta_M(\omega x)=\delta(\omega)\deleted{\unsure{\otimes}}x+\omega\,(\delta_M x)
\qquad\left(\begin{aligned}
&\omega\in\Gamma\bigwedge^*\Omega,\\
&x\in\Gamma\bigwedge^*\Omega\otimes M
\end{aligned}\right)
\]

Pour qu'on ait $\delta_M\delta_M=0$, il f. et s qu'on l'ait \deleted{\unsure{$\delta_M$}} en degré 0, i.e. que le composé
\[
M\xrightarrow{\delta_M^{\unsure{0|\bullet}}}\Omega\otimes M
\overset{\delta_M^1}{\mathbin{\unsure{-|\to}}}\bigwedge^2\Omega\otimes\deleted{\unsure{\wedge}} M
\]
soit nul. Ce composé s'appelle le \uline{tenseur} \uline{courbure}.

La donnée de \inserted{la connexion} $\delta_M^0$ permet de faire opérer $\mathscr T=\underline{\mathrm{Hom}}_{\underline O_X}(\Omega,\underline O_X)$ dans $M$, grâce.
\[
\Theta_X^M(x)=(X\otimes\mathrm{id}_M)(\delta_M^0(x))
\]
On a alors
\[
\begin{array}{r@{\ }l}
\begin{aligned}a)\\[1ex] b)\end{aligned}&
\left\{\begin{aligned}
\Theta_X^M(fx)&=\Theta_X(f)x+f\Theta_X^M(x)\\[1ex]
\Theta_{fX}^M(x)&=f\Theta_X^M(x)
\end{aligned}\right.
\end{array}
\]
D'autre part, si la connexion est sans \deleted{torsion}\inserted{courbure}, alors il est vrai \unsure{qu'on a|que}
\[
c)\quad\left\{\Theta^M_{[X,Y]}=[\Theta_X^{H},\Theta_Y^{H}].\right.
\]

\margin{28}
